% Generated draft: requires mathematical and source review.
\documentclass[11pt]{article}
\usepackage[margin=25mm]{geometry}
\usepackage{fontspec}
\setmainfont{Latin Modern Roman}
\usepackage{amsmath,amssymb,mathtools,unicode-math}
\setmathfont{Latin Modern Math}
\usepackage{enumitem,needspace,xurl,hyperref}
\setlength{\parindent}{0pt}
\setlength{\parskip}{5pt}
\setlength{\emergencystretch}{3em}
\newcommand{\N}{\mathbb{N}}
\newcommand{\Z}{\mathbb{Z}}
\newcommand{\Q}{\mathbb{Q}}
\newcommand{\R}{\mathbb{R}}
\newcommand{\C}{\mathbb{C}}
\newcommand{\F}{\mathbb{F}}
\newcommand{\A}{\mathbb{A}}
\newcommand{\PP}{\mathbb{P}}
\newcommand{\E}{\mathbb{E}}
\newcommand{\eps}{\varepsilon}
\newcommand{\dd}{\,\mathrm{d}}
\begin{document}
\section*{UnsolvedMath ID: 700152}
\textbf{Source problem code:} catalogue entry
\subsection{Definitions and mathematical statement}
For a group \(G\), use the commutator convention \([x,y]=x^{-1}y^{-1}xy\) and define left-normed iterated commutators by
\[
[x,{}_{1}y]=[x,y],\qquad [x,{}_{k+1}y]=[[x,{}_{k}y],y]\quad(k\geq 1).
\]
For an integer \(n\geq 1\), let \(\mathcal E_n\) be the class of groups satisfying the \(n\)th Engel identity:
\[
G\in\mathcal E_n\quad\Longleftrightarrow\quad
(\forall x,y\in G)\;[x,{}_{n}y]=1.
\]
A positive word is a word formed by multiplication of variables without inverse letters. A positive identity is an equation \(u=v\) between such words, required to hold under every assignment of its variables to elements of the group. One may require the words to be nonempty: if an empty word occurs, appending a fresh variable to both sides gives an equivalent identity in groups.

The problem asks whether, for every integer \(n\geq 1\), there exists a family \(\Sigma_n\) of positive identities such that, for every group \(G\),
\[
G\in\mathcal E_n
\quad\Longleftrightarrow\quad
G\models\Sigma_n.
\]
Here \(G\models\Sigma_n\) means that every identity in \(\Sigma_n\) holds for every assignment of its variables in \(G\). The family may depend on \(n\), but must work uniformly for all groups. The source does not specify whether a single identity or a finite family is required; the displayed formulation permits an arbitrary family and leaves any stronger finiteness requirement unresolved.

The supplied statement attributes affirmative results for \(n=1,2,3\) to A. I. Shirshov. The supplied editors' comment attributes an affirmative result for \(n=4\) to G. Traustason. These are source-reported results, not independently verified here.

The editors' comment also records O. Macedo\'{n}ska's observation concerning locally graded \(n\)-Engel groups. A group is locally graded if every nontrivial finitely generated subgroup has a nontrivial finite quotient. A group is locally nilpotent if every finitely generated subgroup is nilpotent. According to the supplied comment, a locally graded \(n\)-Engel group is locally nilpotent by a result of Y. Kim and A. Rhemtulla, and then, by a result of R. G. Burns and Yu. Medvedev, has a normal nilpotent subgroup whose nilpotency class is bounded in terms of \(n\) alone, with quotient of exponent bounded in terms of \(n\) alone. The comment invokes a classical result of Mal'cev to conclude that such groups satisfy a positive law.

The precise scope of this last assertion is unresolved in the supplied record: satisfying a positive law does not by itself state that positive identities characterize the Engel condition. The comment does not explicitly supply an equivalence relative to locally graded groups, a common defining family, or the intended uniformity of the positive law. Its assertion that the result is also true for the locally graded class is therefore retained as a source qualification, without treating it as a precise additional equational characterization.
\subsection{Short English statement}
For each positive integer \(n\), ask whether the groups satisfying \([x,{}_{n}y]=1\) can be characterized, among all groups, by identities involving only products of variables and no inverse letters. The supplied source reports affirmative answers for \(n\leq 4\). Its additional comment about locally graded groups requires clarification about characterization and uniformity. The registry labels the problem open; this label and the reported results have not been independently verified.
\subsection{Sources}
\begin{itemize}
\item Supplied local UnsolvedMath registry record: permanent ID \(700152\), exact source problem code \texttt{catalogue entry}, title \emph{Kourovka Notebook Problem 2.82}. The record labels the problem \emph{open} and was updated on 2026-10-16. This is supplied registry metadata, not an independently verified literature assessment.
\item \emph{The Kourovka Notebook (Unsolved Problems in Group Theory)}, original problem 2.82, second issue (1967); supplied edition reference: twenty-first issue edition (2026), arXiv:1401.0300v48. Proposed by A. I. Shirshov. The supplied collection description names E. I. Khukhro and V. D. Mazurov as editors.
\item A. I. Shirshov, \emph{Algebra i Logika}, volume 2, no. 5 (1963), pp. 5--18, in Russian. Supplied attribution for the cases \(n=1,2,3\); no article title is supplied.
\item G. Traustason, \emph{J. Group Theory}, volume 2, no. 1 (1999), pp. 39--46. Supplied editors' attribution for the case \(n=4\); no article title is supplied.
\item O. Macedo\'{n}ska is credited in the supplied editors' comment with the observation concerning locally graded \(n\)-Engel groups; no separate publication reference is supplied for this observation.
\item Y. Kim and A. Rhemtulla, in \emph{Groups--Korea'94}, de Gruyter, Berlin (1995), pp. 189--197. Cited in the supplied comment for local nilpotence of locally graded \(n\)-Engel groups; no contribution title is supplied.
\item R. G. Burns and Yu. Medvedev, \emph{J. Austral. Math. Soc.}, volume 64 (1998), pp. 92--100. Cited in the supplied comment for the nilpotent-by-bounded-exponent conclusion, with bounds depending only on \(n\); no article title is supplied.
\item A classical result of Mal'cev is invoked in the supplied comment to obtain a positive law. No precise theorem statement or bibliographic reference for that result is supplied.
\end{itemize}
\end{document}
